\chapter{Coordinate Changes and Hand Them Off}\label{ch:handoff}
\section{One writer per boundary}

\begin{headline}
  Use coordination when work crosses repositories, another agent contributes, or ownership
  moves between machines or clients. Each boundary has one writer and a precise packet.
\end{headline}

Start by naming the packet's fields \cite[Dispatch packet; Review]{skill:delegation}:
\begin{itemize}
  \item the deliverable and the evidence;
  \item the writer and its branch or worktree;
  \item forbidden actions;
  \item verification and the handback condition.
\end{itemize}
A precise packet lets a worker finish a bounded contribution without silently inheriting
the parent objective.

\paragraph{Repositories and worktrees} \cite[Start a writing task; Commit and hand
  off]{skill:git}.
\begin{itemize}
  \item Software repositories retain independent histories. A metadata or project-management
        repository owns policy, session records and approved report sources; it is not a
        container for their code history.
  \item Parallel software changes and overlapping report edits use separate worktrees.
  \item A shared checkout can coordinate non-overlapping state paths under one integrator.
  \item Two tasks needing the same report or policy file require serialization or one explicit
        writing owner.
\end{itemize}

\begin{figure}[H]
  \centering
  \begin{tikzpicture}[node distance=6mm and 10mm]
    \node[work,text width=40mm] (c) {Document the input/output contract between the
      layers};
    \node[work,text width=40mm,right=of c] (w) {Each worker gets one exact repository and
      interface};
    \node[check,text width=40mm,right=of w] (i) {The integrator checks the combined evidence
      in dependency order};
    \draw[flow] (c)--(w);
    \draw[flow] (w)--(i);
  \end{tikzpicture}
  \caption{A cross-repository change. One implementation owner owns each physics formula and
    convention. A locally passing helper does not demonstrate that its consumer receives the
    expected units, schema or normalization \cites[Ownership and framework composition;
      Validation order]{skill:implementation}[Dispatch packet; Review]{skill:delegation}.}
  \label{fig:handoff-layers}
\end{figure}

\section{Review a contribution before accepting it}

\begin{headline}
  Review the handback's evidence, not its confidence. Physics and numerical conclusions go
  to a council.
\end{headline}

A useful worker handback contains \cite[Choosing a role; Review]{skill:delegation}:
\begin{itemize}
  \item the actual diff and the checks run;
  \item evidence paths and remaining limits;
  \item a reading route.
\end{itemize}
Independent copy review is read-only feedback to the report writer; domain and
rendered-page review remain separate.

\begin{rulebox}[Mechanical workers]
  A mechanical worker does not decide ambiguous ownership, job control, destructive cleanup,
  or final physics conclusions.
\end{rulebox}

\paragraph{The parent then} \cites[Review]{skill:delegation}[Numerical
  evidence]{skill:implementation}:
\begin{enumerate}
  \item compares the contribution with the original packet and the remaining contract;
  \item resolves findings at their owning layer;
  \item re-runs the affected checks in the relevant environment;
  \item performs the required integration checks once the pieces are coherent;
  \item records deferred or unavailable checks explicitly: they limit the claim the handback can
        support.
\end{enumerate}

\begin{figure}[H]
  \centering
  \begin{tikzpicture}[node distance=5mm and 12mm]
    \node[work,text width=30mm] (brief) {Known fixes first, then one brief};
    \node[check,text width=34mm,right=of brief,yshift=7mm] (m1) {Strong model,\\provider 1};
    \node[check,text width=34mm,right=of brief,yshift=-7mm] (m2) {Strong model,\\provider $n$};
    \node[result,text width=36mm,right=of m1,yshift=-7mm] (v) {Verdicts and findings;
      record who reviewed in which round};
    \draw[flow] (brief.east)--(m1.west);
    \draw[flow] (brief.east)--(m2.west);
    \draw[flow] (m1.east)--(v.west);
    \draw[flow] (m2.east)--(v.west);
    \draw[back] (v.south) to[bend left=30] node[below,lbl]{fix, then round 2}
    (brief.south);
  \end{tikzpicture}
  \caption{The council: one strong model from each provider you can reach reviews the same
    brief in parallel, so a blind spot shared by one model family is less likely to pass
    unseen \cite[Members; Rounds; Record]{skill:council}.}
  \label{fig:handoff-council}
\end{figure}

After two rounds, a finding still open is fixed, or its claim is withdrawn or marked
unsupported. Only a disagreement of judgment that no run can settle becomes a recorded
default or a question to you, and gated actions keep their gates.

\section{Make source durable at the right checkpoint}

\begin{headline}
  Git is the normal source-handoff mechanism across machines. Unique edits are committed
  and pushed before they are left unattended.
\end{headline}

\paragraph{Updating.} Inspect status, branch, worktrees and remotes before an update; pulls
are manual and safe updates fast-forward. Unknown divergence or dirty ownership is
resolved without resetting or overwriting another writer's information \cite[Update
  safely]{skill:git}.

\paragraph{Publishing,} when publication is authorized and classification permits the
destination \cite[Commit and hand off; Material you must not publish]{skill:git}:
\begin{enumerate}
  \item stage only owned paths;
  \item inspect the staged diff;
  \item record checks and limits;
  \item commit coherently;
  \item push to the approved writable remote;
  \item verify the remote commit.
\end{enumerate}
Framework projects have their own fork and contribution rules; consult the repository
instructions before choosing the target.

Scratch is useful for explicitly owned framework checkouts and task worktrees, but it is
not durable merely because it is fast or has space.

\begin{rulebox}[Scratch storage]
  Commit and push unique software edits
  there before leaving them unattended or handing them off. If publication is not authorized
  or fails, record the exact dirty paths and recovery gap in state, and keep the work from
  existing only on scratch \cite[Commit and hand off]{skill:git}.
\end{rulebox}

\begin{table}[H]
  \small
  \begin{tabularx}{\textwidth}{L{0.3\textwidth}Y}
    \toprule
    State & Does not imply \\
    \midrule
    Source checkpoint published & That human report decisions are closed: an authorized
    checkpoint can preserve work while a PDF remains under review. \\
    \rowcolor{ctpSurface0!30}
    PDF delivered & That its source is published. \\
    \bottomrule
  \end{tabularx}
  \caption{Record both states, so the receiving owner knows which evidence is available
    \cites[Commit and hand off]{skill:git}[Review and release]{skill:report-production}.}
  \label{tab:handoff-states}
\end{table}

\section{Give the receiving owner one safe continuation}

\begin{headline}
  Stop the outgoing writer at a coherent checkpoint and hand over one next safe action.
\end{headline}

Give the receiver \cite[Routine continuation and client changes]{skill:recovery}:
\begin{itemize}
  \item the chosen session and exact source state;
  \item current artifacts, active jobs or waits;
  \item remaining checks and the authorization boundary;
  \item the next safe action.
\end{itemize}
The receiver reconciles that record with current Git and runtime evidence before
continuing; Chapter~\ref{ch:daily} distinguishes another display from another writer.

\paragraph{\code{HANDOFF.json}} is useful when a programmatic dependency or cross-harness
recovery needs it \cite[Handoff marker]{skill:jobs}:
\begin{itemize}
  \item it mirrors the Markdown checkpoint atomically;
  \item missing, malformed or held releases do not mean ready;
  \item for a report awaiting human review, the machine status remains \code{in_progress} with
        the report release held.
\end{itemize}

\begin{rulebox}[No speculative release]
  Do not create a release merely because the next step is expected to pass.
\end{rulebox}
